\documentclass[../my_ntust_thesis.tex]{subfiles}
\begin{document}
\iffalse
\chapter{結論}
\fi
\chapter{Conclusion}
\label{sec:conclusions}

This study addresses the transmission latency issues in the O-RAN Open Radio Access Network architecture by proposing an adaptive slots ahead mechanism, with a specific focus on the Option 6 functional split architecture utilizing the nFAPI interface. In a non-ideal packet-switched network environment, traditional static slots ahead configurations cannot effectively handle unpredictable network jitter and server processing delays, easily leading to scheduling timeouts and synchronization failures. To overcome this challenge, this study designs a dynamic delay management algorithm based on the Exponentially Weighted Moving Average (EWMA) model. Through closed-loop feedback control, it continuously tracks and estimates end-to-end network transmission latency characteristics.

By integrating the OpenAirInterface (OAI) open-source software with commercial Pegatron O-RU equipment into an end-to-end physical verification platform, the experimental results demonstrate that the proposed delay management architecture significantly improves system stability. Under simulated network congestion scenarios, this mechanism actively compensates for timing offsets, ensuring that MAC layer scheduling commands accurately fall within the timing window of the physical layer, thereby significantly reducing connection drops caused by missing slot packets within the evaluated range. Compared to traditional solutions that rely on expensive hardware acceleration or strict hardware-level PTP synchronization, the software-based solution proposed in this study demonstrates high deployment flexibility and cost-effectiveness on general-purpose Commercial Off-The-Shelf (COTS) servers, providing a reliable key technological foundation for the commercialization of future 5G/6G heterogeneous networks.

Although the proposed software-based timing control achieves robust performance, the current evaluation focuses primarily on downlink timing control under controlled delay and jitter injection. Future work will extend the proposed method to timing enhancements, multi-cell deployments, and analytical stability bounds for the feedback loop.

\end{document}